\begingroup\small
\setlength\LTleft{0pt}\setlength\LTright{0pt}
\begin{longtable}{@{}p{0.75cm}p{2.35cm}p{2.55cm}p{7.35cm}@{}}
\caption{Representative verified and parametrized domains.}\label{tab:verified}\\
\toprule ID & Domain & Status & Result and caveat\\\midrule\endfirsthead
\toprule ID & Domain & Status & Result and caveat\\\midrule\endhead
VD-01 & Renormalizable Standard Model field content & VERIFIED DOMAIN & All established collider states fit the SU(3)\_C x SU(2)\_L x U(1)\_Y representation pattern; no confirmed extra elementary state. This verifies a low-energy representation pattern, not uniqueness or completeness \\
VD-02 & Higgs sector & VERIFIED DOMAIN & ATLAS combined mass m\_H = 125.11 +/- 0.11 GeV; measured couplings are SM-like within present precision. Self-coupling and total width remain much less precise than the mass \\
VD-03 & Electroweak precision & VERIFIED WITH WATCHLIST & ATLAS m\_W = 80366.5 +/- 15.9 MeV is compatible with electroweak fits. CDF II reports a significantly higher value; this is an inter-experiment tension, not a settled SM failure \\
VD-04 & QCD & VERIFIED DOMAIN & Asymptotic freedom and perturbative factorization succeed at high scales; nonperturbative QCD reproduces extensive hadron and thermodynamic data. Confinement is not derived by a complete analytic theorem in four-dimensional continuum Yang-Mills \\
VD-05 & Flavour and CP violation & VERIFIED DOMAIN WITH LOCAL TENSIONS & A single CKM phase organizes a large body of flavour data. First-row CKM, selected semileptonic observables and hadronic inputs remain precision watchpoints; flavour parameters are fitted, not explained \\
VD-06 & Neutrino sector & EXTENSION REQUIRED & Oscillations establish nonzero mass splittings and mixing; KATRIN gives m\_beta < 0.45 eV (90\% CL). The renormalizable SM without right-handed neutrinos predicts massless neutrinos; the mass origin and Dirac/Majorana character remain open \\
VD-07 & Weak equivalence principle & VERIFIED DOMAIN & MICROSCOPE finds eta(Ti,Pt) = [-1.5 +/- 2.3(stat) +/- 1.5(syst)] x 10\textasciicircum{}-15. The bound is composition- and range-dependent for particular new fields \\
VD-08 & Relativistic binary dynamics & VERIFIED DOMAIN & Multiple strong-field effects agree with GR; quadrupolar radiation damping is tested at roughly 1.3 x 10\textasciicircum{}-4. This does not probe Planck curvature or all strong-field alternatives \\
VD-09 & Gravitational waves & VERIFIED DOMAIN & GWTC-4.0 parameterized tests find no evidence beyond GR; m\_g <= 1.92 x 10\textasciicircum{}-23 eV/c\textasciicircum{}2 at 90\% credibility. Waveform systematics and model flexibility remain part of every precision test \\
VD-10 & Black-hole exterior phenomenology & VERIFIED DOMAIN & EHT and GW observations are consistent with compact objects described by GR black-hole exteriors. Observations do not prove the classical interior or singularity structure \\
VD-11 & Homogeneous cosmological background and CMB & SUCCESSFUL PARAMETRIZATION & Planck finds H0 = 67.4 +/- 0.5 km/s/Mpc, Omega\_m = 0.315 +/- 0.007 and n\_s = 0.965 +/- 0.004 in base LambdaCDM. The fit depends on the model, initial perturbation ansatz and recombination/foreground systematics \\
VD-12 & Large-scale distance ladder and BAO & SUCCESSFUL BASELINE WITH TENSIONS & DESI DR2 BAO is well described by flat LambdaCDM but combined datasets give model- and supernova-sample-dependent preference for evolving dark energy. A preference under an extended parametrization is not an established new component \\
VD-13 & Primordial perturbations & SUCCESSFUL PARAMETRIZATION & CMB data support a near-power-law scalar spectrum; BK18 gives r\_0.05 < 0.036 (95\% CL). No unique inflationary field content or alternative mechanism has been selected \\
VD-14 & Dark matter phenomenology & EMPIRICAL COMPONENT, MICROPHYSICS OPEN & A non-baryonic gravitating component is required in the standard cosmological fit; LZ finds no WIMP excess in 4.2 tonne-years. Null WIMP searches exclude models, not the astrophysical evidence for missing gravitating matter \\
\bottomrule
\end{longtable}
\endgroup
